% latex uft-8
\documentclass[uplatex,a4paper,11pt,oneside,openany]{jsbook}
%
\usepackage[dvipdfmx]{graphicx}
\usepackage[dvipdfmx]{color}
\usepackage{amsmath,amssymb,amsfonts,latexsym}
\usepackage{enumerate}
\usepackage{bm}
\usepackage{graphicx}
\usepackage{ascmac}
\usepackage{setspace}
\usepackage{here}
\usepackage{comment}
\usepackage{bm}
\usepackage{physics}
\usepackage{braket}
\usepackage{fancybox}
\usepackage{fancyhdr}
\usepackage{url}
%\usepackage{jumoline}
\usepackage{listings,jlisting} %日本語のコメントアウトをする場合jlistingが必要

%ここからソースコードの表示に関する設定
\lstset{
	%プログラム言語(複数の言語に対応，C,C++も可)
 	language = Python,
 	%背景色と透過度
 	backgroundcolor={\color[gray]{.95}},
 	%枠外に行った時の自動改行
 	breaklines = true,
 	%自動改行後のインデント量(デフォルトでは20[pt])
 	breakindent = 10pt,
 	%標準の書体
 	%basicstyle = \ttfamily\scriptsize,
  basicstyle=\fontsize{8}{10}\selectfont\ttfamily,
 	%コメントの書体
 	commentstyle = {\itshape \color[cmyk]{1,0.4,1,0}},
 	%関数名等の色の設定
 	classoffset = 0,
 	%キーワード(int, ifなど)の書体
 	keywordstyle = {\bfseries \color[cmyk]{0,1,0,0}},
 	%表示する文字の書体
 	stringstyle = {\ttfamily \color[rgb]{0,0,1}},
 	%枠 "t"は上に線を記載, "T"は上に二重線を記載
	%他オプション：leftline，topline，bottomline，lines，single，shadowbox
 	frame = TBrl,
 	%frameまでの間隔(行番号とプログラムの間)
 	framesep = 5pt,
 	%行番号の位置
 	numbers = left,
	%行番号の間隔
 	stepnumber = 1,
	%行番号の書体
 	numberstyle = \tiny,
	%タブの大きさ
 	tabsize = 4,
 	%キャプションの場所("tb"ならば上下両方に記載)
 	captionpos = t
}
%ここまでソースコードの表示に関する設定
%ここからソースコードの表示に関する設定
%\lstset{
%  basicstyle={\ttfamily},
%  identifierstyle={\small},
%  commentstyle={\smallitshape},
%  keywordstyle={\small\bfseries},
%  ndkeywordstyle={\small},
%  stringstyle={\small\ttfamily},
%  frame={tb},
%  breaklines=true,
%  columns=[l]{fullflexible},
%  numbers=left,
%  xrightmargin=0zw,
%  xleftmargin=3zw,
%  numberstyle={\scriptsize},
%  stepnumber=1,
%  numbersep=1zw,
%  lineskip=-0.5ex
%}
%ここまでソースコードの表示に関する設定
\makeatletter
\def\ps@plainfoot{%
  \let\@mkboth\@gobbletwo
  \let\@oddhead\@empty
  \def\@oddfoot{\normalfont\hfil-- \thepage\ --\hfil}%
  \let\@evenhead\@empty
  \let\@evenfoot\@oddfoot}
  \let\ps@plain\ps@plainfoot
%\@openrighttrue
\renewcommand{\chapter}{%
  \if@openright\cleardoublepage\else\clearpage\fi
  \global\@topnum\z@
  \secdef\@chapter\@schapter}
\makeatother
%
\newcommand{\maru}[1]{{\ooalign{%
\hfil\hbox{$\bigcirc$}\hfil\crcr%
\hfil\hbox{#1}\hfil}}}
%
\setlength{\textwidth}{\fullwidth}
\setlength{\textheight}{40\baselineskip}
\addtolength{\textheight}{\topskip}
\setlength{\voffset}{-0.55in}
%
%\pagestyle{footnombre}
%
\begin{document}
% START DOCUMENT
%
%\pagestyle{fancy}
%\fancyhead{}
%\fancyhead[RE]{\rightmark}
%\fancyhead[LO]{\leftmark}
%\fancyhead[LE,RO]{\thepage}
%\fancyfoot{}
%
% COVER
\begin{center}
  \huge \_ \par
  \vspace{25mm}
  \LARGE 2022年度 \par
  \vspace{15mm}
  \huge ダーツ・ゲーム機の製作 \par
  \vspace{100mm}
  \Large \today \par
  \vspace{15mm}
  \Large S.Matoike \par
  \vspace{10mm}
  \Large \par
  \vspace{10mm}
\end{center}
\thispagestyle{empty}
\clearpage
\addtocounter{page}{-1}
\newpage
\setcounter{tocdepth}{3}
%
\tableofcontents
%
\chapter{はじめに}
%
\begin{align*}
合計：\;\;& \sum_{i=1}^n x_i=x_1+x_2+\cdots +x_n\\
平均：\;\;&\bar{x} = \frac{1}{n}\sum_{i=1}^n x_i = \frac{1}{n}(x_1+x_2+\cdots+x_n)\\\\
偏差：\;\;&x_i-\bar{x}\\\\
分散：\sigma^2 &= \frac{1}{n}\sum_{i=1}^n(x_i-\bar{x})^2=\frac{1}{n}\{ (x_1-\bar{x})^2 + (x_2-\bar{x})^2+\dots+(x_n-\bar{x})^2  \}\\
標準偏差：\sigma&=\sqrt{\sigma^2}\\\\
%標準偏差：&\sigma=\sqrt{\sigma^2}=\sqrt{\frac{1}{n}\sum_{i=1}^n(x_i-\bar{x})^2}\\\\
基準化変量：&z_i=\frac{x_i-\bar{x}}{\sigma}\\\\
偏差値：&\;\;50+10\times z_i=50+10\times\frac{x_i-\bar{x}}{\sigma}
\end{align*}

\chapter{おわりに}
%
\section*{謝辞}
\addcontentsline{toc}{chapter}{謝辞}
%
\appendix
%
\begin{thebibliography}{99}
  \bibitem{1}
\end{thebibliography}
%
% END DOCUMENT
\end{document}
%
